\subsection{Data-Driven Design Hypotheses}
\label{sec:hypotheses}

\todo{Draft from the outline.}
\begin{itemize}
  \item \textbf{Hypothesis $H_1$ (continuous log-regression for $\Dch$):} modeling $Y = \log_{10} \Dch$ directly minimizes relative error, achieving sub-percent MARE across all $6.7$ decades without requiring a complex classify-then-regress pipeline.
  \item \textbf{Hypothesis $H_2$ (turning-point gradient protection):} combining Huber loss ($\delta = 1.0$) with smooth GELU activations prevents gradient saturation at $dM/d\rhoc = 0$, reducing the $\Mmax$ boundary error by $>50\%$ versus standard MSE/ReLU networks.
  \item \textbf{Hypothesis $H_3$ (ResNet grid smoothing):} identity skip-connections ($h^{(l+1)} = h^{(l)} + F(h^{(l)})$) allow the network to model continuous $\log\rhoc$ variation while smoothing over discrete $(\beta, \lambda)$ grid steps without residual ripples.
\end{itemize}
